\documentclass[11pt,a4paper]{article}

% ------------------------------------------------
% Paquetes
% ------------------------------------------------
\usepackage[utf8]{inputenc}
\usepackage[T1]{fontenc}
\usepackage[spanish,es-nodecimaldot]{babel}

\usepackage{amsmath}
\usepackage{amssymb}
\usepackage{array}
\usepackage{booktabs}
\usepackage{geometry}
\usepackage{enumitem}
\usepackage{listings}
\usepackage{xcolor}
\usepackage{float}
\usepackage{tocloft}

% ------------------------------------------------
% Configuración de página
% ------------------------------------------------
\geometry{
    margin=2.5cm
}

% ------------------------------------------------
% Configuración de código
% ------------------------------------------------
\lstset{
    basicstyle=\ttfamily\small,
    columns=fullflexible,
    keepspaces=true,
    frame=single,
    breaklines=true
}

% ------------------------------------------------
% Documento
% ------------------------------------------------
\begin{document}

\begin{center}
    {\LARGE \textbf{Resumen ampliado --- Capítulo 1}}\\[4pt]
    {\Large \textit{Introducción al diseño e implementación de compiladores}}\\[4pt]
    \textit{Basado en el Capítulo 1 de Aho, Lam, Sethi y Ullman, \textbf{Compiladores: principios,
    técnicas y herramientas} (``Libro del Dragón'')}\\[2pt]
    \small Material de repaso --- Cátedra de Compiladores
\end{center}

\vspace{0.4cm}

\tableofcontents

\bigskip

\section{Introducción}

Los lenguajes de programación son la forma en que las personas describen cálculos a las
máquinas. Sin embargo, ningún programa fuente puede ejecutarse directamente: antes debe
\textbf{traducirse} a un formato que la computadora pueda procesar. Los sistemas de software
que realizan esta traducción se llaman \textbf{compiladores}.

El estudio de los compiladores no es solamente el estudio de una herramienta puntual: las
ideas y técnicas que se utilizan para construirlos (autómatas finitos, expresiones regulares,
gramáticas libres de contexto, análisis de flujo de datos, optimización) reaparecen en muchas
otras áreas de las ciencias de la computación. Por esta razón, comprender la estructura general
de un compilador es la base indispensable para todo lo que veremos después en la materia,
en particular para las herramientas \textbf{Flex} (generador de analizadores léxicos) y
\textbf{Bison} (generador de analizadores sintácticos), que automatizan justamente las dos
primeras fases que se describen en este capítulo.

\section{Procesadores de lenguaje}

En su forma más simple, un \textbf{compilador} es un programa que lee un programa escrito en
un \textbf{lenguaje fuente} y lo traduce a un programa equivalente escrito en un
\textbf{lenguaje destino}. Una función esencial del compilador es, además, reportar los
errores del programa fuente que detecte durante la traducción.

\[
\text{programa fuente} \longrightarrow \boxed{\text{Compilador}} \longrightarrow \text{programa destino}
\]

Si el programa destino es un programa ejecutable en lenguaje máquina, el usuario puede
ejecutarlo directamente para procesar entradas y producir salidas:

\[
\text{entrada} \longrightarrow \boxed{\text{Programa destino}} \longrightarrow \text{salida}
\]

\subsection{Compiladores vs. intérpretes}

Un \textbf{intérprete} es otro tipo de procesador de lenguaje muy común. En vez de producir un
programa destino como traducción, el intérprete da la apariencia de ejecutar directamente las
operaciones especificadas en el programa fuente, usando las entradas suministradas por el
usuario:

\[
\text{programa fuente, entrada} \longrightarrow \boxed{\text{Intérprete}} \longrightarrow \text{salida}
\]

En términos de eficiencia y diagnóstico de errores existe un compromiso:

\begin{itemize}
    \item El programa destino en lenguaje máquina que produce un \textbf{compilador} suele ser
    \textbf{más rápido} que un intérprete al mapear entradas a salidas, porque la traducción se
    realiza una sola vez.
    \item Un \textbf{intérprete} suele ofrecer \textbf{mejores diagnósticos de error}, ya que
    ejecuta el programa fuente instrucción por instrucción y puede reportar el contexto exacto en
    el que ocurre un fallo.
\end{itemize}

\textbf{Procesadores híbridos.} Muchos lenguajes modernos combinan ambas estrategias. El caso
clásico es Java: el programa fuente se compila primero a un formato intermedio llamado
\textit{bytecode}, y luego una máquina virtual interpreta esos bytecodes. Esto permite que el
mismo bytecode compilado en una máquina se ejecute en otra distinta (incluso a través de una
red). Para acelerar aún más la ejecución, los compiladores \textit{just-in-time} (JIT) traducen
los bytecodes a lenguaje máquina justo antes de ejecutarlos.

\[
\begin{array}{c}
\text{programa fuente} \to \boxed{\text{Traductor}} \to \text{programa intermedio} \\[4pt]
\downarrow \text{entrada} \\[2pt]
\boxed{\text{Máquina virtual}} \to \text{salida}
\end{array}
\]

\subsection{El sistema completo de procesamiento de lenguaje}

Para producir un ejecutable final, generalmente se requieren varios programas además del
compilador propiamente dicho:

\begin{enumerate}
    \item Un \textbf{preprocesador} recolecta el programa fuente (que puede estar dividido en
    varios módulos/archivos) y expande \textit{macros}, es decir, fragmentos de código
    abreviados de uso frecuente, en instrucciones del lenguaje fuente.
    \item El \textbf{compilador} traduce el programa fuente (ya preprocesado) a un programa
    destino, que muchas veces es código en \textbf{lenguaje ensamblador}, porque resulta más
    fácil de generar y de depurar que el código máquina directo.
    \item Un \textbf{ensamblador} procesa el código ensamblador y produce \textbf{código máquina
    relocalizable}.
    \item Un \textbf{enlazador (linker)} resuelve las referencias a direcciones de memoria
    externas, es decir, cuando código de un archivo hace referencia a algo definido en otro
    archivo, y combina el código relocalizable con archivos objeto y bibliotecas.
    \item Un \textbf{cargador (loader)} reúne todos los archivos objeto ejecutables en memoria
    para su ejecución.
\end{enumerate}

\[
\begin{array}{c}
\text{fuente} \to \boxed{\text{Preprocesador}} \to \text{fuente modificado} \to
\boxed{\text{Compilador}} \to \text{ensamblador} \\[6pt]
\downarrow \\[6pt]
\boxed{\text{Ensamblador}} \to \text{código máquina relocalizable} \\[6pt]
\downarrow \\[6pt]
\boxed{\text{Enlazador/Cargador}} \to \text{código máquina destino}
\end{array}
\]

Esta separación en programas independientes es una primera manifestación de una idea que
reaparecerá constantemente en la materia: dividir un problema grande y complejo (traducir un
lenguaje) en \textbf{etapas más pequeñas y manejables}, cada una con una responsabilidad
claramente delimitada.

\section{La estructura interna de un compilador}

Si ``abrimos la caja'' del compilador, encontramos que la traducción se realiza en dos grandes
partes:

\begin{itemize}
    \item El \textbf{análisis} (\textit{front-end}) descompone el programa fuente en sus
    componentes, les impone una estructura gramatical y construye a partir de ellos una
    \textbf{representación intermedia}. Si detecta que el programa está mal formado
    sintácticamente o que su semántica es inconsistente, debe emitir mensajes de error útiles.
    El análisis también recolecta información sobre el programa en la \textbf{tabla de
    símbolos}.
    \item La \textbf{síntesis} (\textit{back-end}) construye el programa destino a partir de la
    representación intermedia y de la información almacenada en la tabla de símbolos.
\end{itemize}

En nuestro caso, el lenguaje destino de todo el proceso será \textbf{C}, de modo que el
compilador que construiremos terminará generando un programa en C, el cual luego será
compilado por un compilador de C ya existente.

De forma general:

\[
\text{Programa fuente} \rightarrow \text{Análisis} \rightarrow \text{Representación intermedia}
\rightarrow \text{Optimización} \rightarrow \text{Código destino}
\]

Analizado con más detalle, el proceso de compilación opera como una \textbf{secuencia de
fases}, cada una de las cuales transforma una representación del programa fuente en otra:

\[
\begin{array}{c}
\text{flujo de caracteres} \\
\downarrow \\
\boxed{\text{Analizador léxico}} \\
\downarrow \\
\text{flujo de tokens} \\
\downarrow \\
\boxed{\text{Analizador sintáctico}} \\
\downarrow \\
\text{árbol sintáctico} \\
\downarrow \\
\boxed{\text{Analizador semántico}} \\
\downarrow \\
\text{árbol sintáctico (anotado)} \\
\downarrow \\
\boxed{\text{Generador de código intermedio}} \\
\downarrow \\
\text{representación intermedia} \\
\downarrow \\
\boxed{\text{Optimizador de código}} \\
\downarrow \\
\text{representación intermedia optimizada} \\
\downarrow \\
\boxed{\text{Generador de código}} \\
\downarrow \\
\text{código máquina destino}
\end{array}
\]

La \textbf{tabla de símbolos} se utiliza de manera transversal en todas las fases, ya que
almacena información sobre todo el programa fuente. La fase de \textbf{optimización de código
independiente de la máquina}, entre el front-end y el back-end, es \textbf{opcional}: su
objetivo es transformar la representación intermedia para que el back-end pueda producir mejor
código destino.

En la práctica, varias fases pueden agruparse y las representaciones intermedias entre fases
agrupadas no necesitan construirse explícitamente; volveremos sobre esta idea en la sección
\ref{sec:fases-pasadas} (\textit{Fases y pasadas}).

Como ejemplo guía utilizaremos el siguiente programa fuente, igual que en el resto de la
cátedra:

\begin{lstlisting}
int x = 10 + 20;
\end{lstlisting}

\subsection{Análisis léxico}
\label{sec:lexico}

El análisis léxico, también llamado \textbf{escaneo}, constituye la primera fase del
compilador. El analizador léxico lee el flujo de caracteres del programa fuente y los agrupa en
secuencias significativas llamadas \textbf{lexemas}. Para cada lexema, produce como salida un
\textbf{token} de la forma:

\[
\langle \text{nombre-token}, \text{valor-atributo} \rangle
\]

que se pasa a la fase siguiente, el análisis sintáctico. El \textbf{nombre-token} es un símbolo
abstracto utilizado durante el análisis sintáctico (por ejemplo, \texttt{id}, \texttt{+}), y el
\textbf{valor-atributo} suele apuntar a una entrada en la \textbf{tabla de símbolos} con
información adicional sobre ese token (por ejemplo, el nombre y el tipo de un identificador).

De esta manera, el análisis léxico separa la \textbf{representación concreta} escrita por el
programador de la \textbf{categoría} que interesa a las fases posteriores, y el analizador
léxico \textbf{ignora los espacios en blanco} que separan a los lexemas.

Por lo tanto, el análisis léxico responde principalmente a la pregunta:

\begin{quote}
\textit{¿Qué unidades léxicas aparecen en la secuencia de caracteres del programa fuente?}
\end{quote}

El analizador léxico \textbf{no} determina si esas unidades pueden combinarse correctamente:
esa responsabilidad corresponde al analizador sintáctico (sección \ref{sec:sintactico}).

Para nuestro ejemplo:

\begin{lstlisting}
int x = 10 + 20;
\end{lstlisting}

podemos identificar los siguientes lexemas:

\begin{lstlisting}
int
x
=
10
+
20
;
\end{lstlisting}

Cada uno de estos lexemas pertenece a una categoría léxica (token):

\begin{lstlisting}
INT
IDENTIFIER
ASSIGN
INTEGER
PLUS
INTEGER
SEMICOLON
\end{lstlisting}

\paragraph{Relevancia para Flex.} Esta fase es exactamente la que automatiza \textbf{Flex}: a
partir de una especificación de patrones (expresiones regulares) para cada categoría léxica,
Flex genera automáticamente el código en C de una función (típicamente \texttt{yylex()}) capaz
de leer el flujo de caracteres, agruparlo en lexemas y devolver, uno por uno, los tokens
correspondientes. Retomaremos esta conexión en la sección \ref{sec:flex-bison}.

\subsection{Lexemas, tokens y patrones}
\label{sec:lexemas-tokens-patrones}

Es fundamental distinguir con precisión estos tres conceptos, ya que son la base conceptual de
cualquier especificación léxica (y, por lo tanto, de cualquier archivo \texttt{.l} de Flex):

\begin{itemize}
    \item Un \textbf{lexema} es una secuencia concreta de caracteres del programa fuente que
    forma una unidad, por ejemplo \texttt{contador}, \texttt{25} o \texttt{=}.
    \item Un \textbf{token} es la categoría léxica que se asigna a un lexema, representada como
    un par $\langle \text{nombre-token}, \text{valor-atributo} \rangle$. El
    \textbf{valor-atributo} sólo es necesario cuando existen varios lexemas posibles para el
    mismo token (por ejemplo, distintos identificadores); tokens como \texttt{=} no lo
    necesitan.
    \item Un \textbf{patrón} es la \textbf{regla} —típicamente una expresión regular— que
    describe qué secuencias de caracteres pueden formar los lexemas de un determinado token.
\end{itemize}

Para nuestro lenguaje de ejemplo, la especificación léxica inicial es:

\begin{table}[H]
\centering
\begin{tabular}{>{\ttfamily}c c c}
\toprule
\textbf{Token} & \textbf{Patrón (expresión regular)} & \textbf{Ejemplo de lexema} \\
\midrule
INT &
\texttt{int} &
\texttt{int} \\

IDENTIFIER &
\texttt{[a-zA-Z\_][a-zA-Z0-9\_]*} &
\texttt{x}, \texttt{contador} \\

INTEGER &
\texttt{[0-9]+} &
\texttt{10}, \texttt{250} \\

ASSIGN &
\texttt{=} &
\texttt{=} \\

PLUS &
\texttt{+} &
\texttt{+} \\

SEMICOLON &
\texttt{;} &
\texttt{;} \\
\bottomrule
\end{tabular}
\caption{Especificación léxica inicial del lenguaje}
\label{tab:tokens}
\end{table}

Nótese que la columna ``patrón'' de la tabla \ref{tab:tokens} \textbf{es, literalmente}, lo que
se escribiría en la sección de definiciones de tokens de un archivo de especificación de Flex:
cada regla de Flex es, en esencia, un par (\textit{patrón} en forma de expresión regular,
\textit{acción} en C que construye y devuelve el token correspondiente).

Para el programa completo:

\begin{lstlisting}
int x = 10 + 20;
\end{lstlisting}

el analizador léxico produce la secuencia de tokens:

\begin{lstlisting}
INT
IDENTIFIER(x)
ASSIGN
INTEGER(10)
PLUS
INTEGER(20)
SEMICOLON
\end{lstlisting}

Los tokens que poseen un valor asociado (como los identificadores y los literales enteros)
transportan esa información adicional, por ejemplo \texttt{IDENTIFIER(x)} o
\texttt{INTEGER(10)}. El identificador también puede representarse mediante una referencia a
una entrada de la tabla de símbolos, siguiendo la notación clásica del Libro del Dragón:

\[
(\,\mathbf{id},\ 1\,)
\]

donde $\mathbf{id}$ representa la categoría del token (identificador) y $1$ identifica la
entrada correspondiente en la tabla de símbolos (sección \ref{sec:tabla-simbolos}).

\subsection{Análisis sintáctico}
\label{sec:sintactico}

El análisis sintáctico, o \textit{parsing}, es la segunda fase del compilador. El analizador
sintáctico (\textit{parser}) recibe la secuencia de tokens producida por el analizador léxico y
\textbf{determina si esa secuencia puede derivarse} de acuerdo con la gramática del lenguaje,
construyendo además una representación intermedia en forma de árbol —típicamente un
\textbf{árbol sintáctico}— que describe la estructura gramatical del flujo de tokens.

Mientras que el análisis léxico determina \textbf{qué} unidades aparecen en el programa, el
análisis sintáctico determina \textbf{cómo} pueden combinarse esas unidades. Por lo tanto,
responde a la pregunta:

\begin{quote}
\textit{¿Cómo se combinan los tokens para formar construcciones válidas del lenguaje?}
\end{quote}

Por ejemplo, la secuencia:

\begin{lstlisting}
INT IDENTIFIER ASSIGN INTEGER PLUS INTEGER SEMICOLON
\end{lstlisting}

corresponde a una declaración válida (\texttt{int x = 10 + 20;}), mientras que:

\begin{lstlisting}
INT IDENTIFIER ASSIGN PLUS INTEGER SEMICOLON
\end{lstlisting}

contiene tokens individualmente válidos, pero cuya organización \textbf{no} corresponde a
ninguna construcción permitida por la gramática: es un \textbf{error sintáctico}.

En un árbol sintáctico, cada \textbf{nodo interior} representa una operación, y los
\textbf{hijos} de ese nodo representan sus argumentos. Esta representación en árbol es la que
permite, por ejemplo, expresar de forma no ambigua la \textbf{precedencia de operadores}: para
\texttt{10 + 20 * 3}, el árbol debe reflejar que la multiplicación se agrupa primero, de modo
que la expresión se interprete como $10 + (20 \times 3)$ y no como $(10+20)\times 3$.

\paragraph{Relevancia para Bison.} Esta es exactamente la fase que automatiza \textbf{Bison}: a
partir de una \textbf{gramática libre de contexto} (una colección de reglas de producción como
las de la sección \ref{sec:especificacion-sintactica}), Bison genera un analizador LALR capaz
de reconocer si una secuencia de tokens es derivable desde el símbolo inicial de la gramática,
y de ejecutar \textbf{acciones semánticas} asociadas a cada regla —típicamente, la construcción
de nodos de un árbol de sintaxis abstracta—.

\subsection{Especificación sintáctica}
\label{sec:especificacion-sintactica}

La sintaxis de un lenguaje de programación se especifica formalmente mediante una
\textbf{gramática libre de contexto} (GLC). Una gramática de este tipo consiste en un conjunto
de \textbf{símbolos no terminales} (categorías sintácticas, como $\langle expresion \rangle$),
\textbf{símbolos terminales} (los tokens producidos por el analizador léxico) y
\textbf{producciones} que indican cómo un no terminal puede reescribirse en una secuencia de
terminales y no terminales.

Para el lenguaje propuesto, una gramática inicial es:

\begin{align*}
\langle programa \rangle
    &\rightarrow \langle declaracion \rangle \\[4pt]
\langle declaracion \rangle
    &\rightarrow
    \langle tipo \rangle
    \ \texttt{IDENTIFIER}
    \ \texttt{ASSIGN}
    \ \langle expresion \rangle
    \ \texttt{SEMICOLON} \\[4pt]
\langle tipo \rangle
    &\rightarrow \texttt{INT} \\[4pt]
\langle expresion \rangle
    &\rightarrow
    \langle expresion \rangle
    \ \texttt{PLUS}
    \ \langle termino \rangle \\[4pt]
\langle expresion \rangle
    &\rightarrow \langle termino \rangle \\[4pt]
\langle termino \rangle
    &\rightarrow \texttt{INTEGER}
\end{align*}

Esta gramática establece qué secuencias de tokens forman declaraciones y expresiones válidas.
Nótese que la producción recursiva de $\langle expresion \rangle$ (que se refiere a sí misma
del lado izquierdo, con \texttt{PLUS} como operador entre una expresión y un término) es lo que
permite generar expresiones con una cantidad arbitraria de sumandos, y su forma
\textbf{recursiva a izquierda} determina que \texttt{+} se agrupe de izquierda a derecha; el
tratamiento riguroso de este tipo de ambigüedades y de la precedencia de operadores mediante
gramáticas es, precisamente, uno de los temas centrales que veremos en la clase de gramáticas.

Para el programa \texttt{int x = 10 + 20;}, el analizador sintáctico puede construir un árbol
de derivación equivalente a:

\begin{lstlisting}
declaracion
|-- tipo -> INT
|-- IDENTIFIER -> x
|-- ASSIGN -> =
|-- expresion
|   |-- expresion
|   |   `-- termino
|   |       `-- INTEGER -> 10
|   |-- PLUS
|   `-- termino
|       `-- INTEGER -> 20
`-- SEMICOLON
\end{lstlisting}

Esta estructura constituye una representación de la organización sintáctica del programa; es lo
que en Bison se conoce como \textbf{árbol de derivación} o \textit{parse tree}, y suele ser más
detallado (y más cercano a la gramática) de lo estrictamente necesario para las fases
posteriores.

\subsection{Árbol de sintaxis abstracta}
\label{sec:ast}

A partir del análisis sintáctico se puede construir un \textbf{árbol de sintaxis abstracta}
(\textit{Abstract Syntax Tree}, AST). A diferencia del árbol de derivación completo, el AST
\textbf{elimina detalles sintácticos} que no son necesarios para las fases posteriores (como
paréntesis, palabras clave redundantes o nodos intermedios de la gramática) y conserva
principalmente la \textbf{estructura semánticamente relevante} del programa.

Para \texttt{int x = 10 + 20;}, un AST puede representarse como:

\begin{lstlisting}
Declaration
|-- type: Int
|-- name: x
`-- value: BinaryExpression
    |-- left: IntegerLiteral(10)
    |-- operator: +
    `-- right: IntegerLiteral(20)
\end{lstlisting}

El AST permite que las fases siguientes (análisis semántico, generación de código intermedio)
trabajen sobre una representación estructurada del programa, en lugar de procesar directamente
los caracteres del código fuente o los detalles superfluos del árbol de derivación.

Además, la estructura del árbol permite representar propiedades sintácticas como la
\textbf{precedencia de operadores}. Por ejemplo:

\begin{lstlisting}
10 + 20 * 3
\end{lstlisting}

debe interpretarse como $10 + (20 \times 3)$ y no como $(10+20)\times 3$. El AST representa esta
diferencia mediante la \textbf{posición de los nodos dentro del árbol}: el nodo \texttt{*}
queda anidado como hijo del nodo \texttt{+}, y no al revés.

\paragraph{Relación con Bison.} En una especificación de Bison, la construcción del AST se
realiza típicamente mediante \textbf{acciones semánticas} escritas en C, asociadas a cada regla
de la gramática (por ejemplo, con el pseudo-código \texttt{\$\$ = nuevoNodoBinario(\$1, PLUS,
\$3);} dentro de la regla que reconoce $\langle expresion \rangle \rightarrow \langle expresion
\rangle\ \texttt{PLUS}\ \langle termino \rangle$). Esta técnica se conoce como
\textbf{traducción dirigida por la sintaxis}, y es el mecanismo central que conecta la gramática
formal con la construcción efectiva del AST en C.

\subsection{Análisis semántico}
\label{sec:semantico}

El hecho de que un programa sea \textbf{sintácticamente correcto} no implica que sea
\textbf{semánticamente válido}. El analizador semántico utiliza el árbol sintáctico (o el AST)
y la información de la tabla de símbolos para comprobar restricciones que no pueden expresarse
adecuadamente mediante una gramática libre de contexto —por ejemplo, que un identificador haya
sido declarado antes de usarse, o que los tipos de una operación sean compatibles—.

Una parte central del análisis semántico es la \textbf{comprobación (verificación) de tipos}:
el compilador verifica que cada operador tenga operandos del tipo esperado. Cuando el lenguaje
lo permite, pueden aplicarse \textbf{coerciones}, es decir, conversiones automáticas de tipo
(por ejemplo, de \texttt{int} a \texttt{float}); volveremos sobre esto con un ejemplo completo
en la sección \ref{sec:ejemplo-completo}.

Por ejemplo:

\begin{lstlisting}
int x = "hola";
\end{lstlisting}

puede tener una estructura \textbf{sintácticamente válida}, ya que respeta la forma general
\textit{tipo, identificador, asignación, expresión} exigida por la gramática, pero es
\textbf{semánticamente incorrecta}, porque una cadena no puede asignarse a una variable de tipo
\texttt{int}.

Para una declaración genérica \texttt{int x = e;} podemos establecer las siguientes reglas
semánticas:

\begin{enumerate}
    \item $x$ debe ser un identificador válido.
    \item $x$ no debe haber sido declarado previamente en el mismo ámbito.
    \item La expresión $e$ debe producir un valor de tipo \texttt{int} (o un tipo que pueda
    coercionarse a \texttt{int}).
    \item El valor producido por $e$ se asigna a $x$.
    \item $x$ queda registrada en la tabla de símbolos como una variable de tipo \texttt{int}.
\end{enumerate}

Para la expresión \texttt{10 + 20} se puede determinar:

\[
10 : \texttt{int} \qquad 20 : \texttt{int} \qquad 10 + 20 : \texttt{int}
\]

Por lo tanto, \texttt{int x = 10 + 20;} es semánticamente válida.

\subsection{Tabla de símbolos}
\label{sec:tabla-simbolos}

La \textbf{tabla de símbolos} es una estructura de datos utilizada por distintas fases del
compilador para almacenar información sobre los nombres (identificadores) que aparecen en el
programa: variables, funciones, tipos, etc. \textbf{No} constituye necesariamente una fase
independiente, sino un \textbf{servicio transversal} utilizado sobre todo durante el análisis
semántico y, posteriormente, durante la generación de código.

Entre los atributos que puede almacenar una entrada se encuentran:

\begin{itemize}
    \item nombre;
    \item tipo;
    \item ámbito (\textit{scope});
    \item ubicación en memoria;
    \item para funciones/procedimientos: número y tipo de parámetros, forma de pasaje (por
    valor o por referencia) y tipo de retorno;
    \item información adicional necesaria para la generación de código.
\end{itemize}

La estructura de datos debe diseñarse para permitir \textbf{buscar, insertar y actualizar} cada
entrada con rapidez (típicamente mediante una tabla hash o un árbol balanceado).

Para \texttt{int x = 10 + 20;} obtenemos:

\begin{table}[H]
\centering
\begin{tabular}{ccc}
\toprule
\textbf{Nombre} & \textbf{Tipo} & \textbf{Ámbito} \\
\midrule
\texttt{x} & \texttt{int} & local \\
\bottomrule
\end{tabular}
\caption{Ejemplo de tabla de símbolos}
\label{tab:symbol-table}
\end{table}

Cuando el analizador léxico encuentra el lexema \texttt{x}, produce un token asociado a una
entrada de esta tabla: \texttt{IDENTIFIER -> (id, 1)}. Las fases posteriores utilizan esta
referencia para obtener toda la información disponible sobre el identificador, en lugar de
volver a manipular directamente la cadena de caracteres \texttt{x}.

\subsection{Manejo de errores}
\label{sec:errores}

El compilador también debe \textbf{detectar y comunicar} los errores que encuentre durante el
procesamiento del programa. Los errores pueden aparecer en cualquier fase:

\begin{itemize}
    \item un carácter que no pertenece al lenguaje produce un \textbf{error léxico} (por
    ejemplo, un símbolo no reconocido por ningún patrón);
    \item una secuencia de tokens que no puede derivarse mediante la gramática produce un
    \textbf{error sintáctico};
    \item una operación incompatible con los tipos involucrados produce un \textbf{error
    semántico}.
\end{itemize}

El manejo de errores es una actividad \textbf{transversal} a todo el proceso de compilación. El
objetivo no es solamente detectar que existe un error, sino proporcionar información suficiente
para que el programador pueda identificar su \textbf{ubicación} (por ejemplo, línea y columna) y
su \textbf{naturaleza}. Tanto Flex como Bison ofrecen mecanismos básicos para esto: Flex puede
reportar un carácter no reconocido por ningún patrón, y Bison invoca una rutina de manejo de
error (habitualmente \texttt{yyerror}) cuando la secuencia de tokens no es derivable desde la
gramática.

\subsection{Generación de código intermedio}
\label{sec:codigo-intermedio}

Después del análisis, el compilador puede construir una o más \textbf{representaciones
intermedias} (\textit{Intermediate Representation}, IR). Los árboles sintácticos (y el AST) son
en sí mismos una forma de representación intermedia, utilizada durante el análisis sintáctico y
semántico.

Muchos compiladores generan además una representación intermedia de \textbf{más bajo nivel},
similar a un lenguaje ensamblador abstracto para una máquina hipotética, que debe cumplir dos
propiedades importantes: ser \textbf{fácil de producir} (a partir del AST) y \textbf{fácil de
traducir} a la máquina destino. Una forma habitual de IR es el \textbf{código de tres
direcciones}: una secuencia de instrucciones simples, con a lo sumo un operador y tres
operandos por instrucción (cada operando puede actuar como un registro).

Para \texttt{int x = 10 + 20;}, una representación intermedia sencilla podría ser:

\begin{lstlisting}
t1 = 10 + 20
x = t1
\end{lstlisting}

La representación intermedia permite \textbf{separar el lenguaje fuente del lenguaje/arquitectura
destino}: en lugar de generar directamente código específico de una máquina, el compilador
representa las operaciones del programa de forma independiente de la plataforma. Esto es
justamente lo que aprovecharemos al usar \textbf{C como lenguaje destino} (sección
\ref{sec:c-destino}): nuestro back-end sólo necesita traducir la IR a C, y delegamos en el
compilador de C el trabajo de generar el código específico de cada arquitectura.

\subsection{Optimización de código}
\label{sec:optimizacion}

Una vez construida la representación intermedia, pueden aplicarse transformaciones que
\textbf{preserven el significado} del programa mientras mejoran alguna de sus propiedades:
tiempo de ejecución, tamaño del código generado o consumo de energía, entre otras.

Por ejemplo:

\begin{lstlisting}
t1 = 10 + 20
x = t1
\end{lstlisting}

puede transformarse en:

\begin{lstlisting}
x = 30
\end{lstlisting}

Esta transformación, en la que una expresión constante se evalúa en tiempo de compilación, se
conoce como \textbf{plegado de constantes} (\textit{constant folding}). Las optimizaciones
pueden aplicarse sobre expresiones, instrucciones, flujo de control o el uso de valores
temporales, entre muchos otros aspectos.

Es fundamental notar que la optimización \textbf{no modifica el significado observable} del
programa: su único objetivo es obtener una representación equivalente que se ejecute de manera
más eficiente. Esta fase es \textbf{opcional} y su alcance varía mucho entre compiladores.

\subsection{Generación de código destino}
\label{sec:generacion-codigo}

La última fase principal del compilador transforma la representación intermedia (optimizada o
no) en el \textbf{lenguaje destino}. En un compilador tradicional, esto suele ser lenguaje
ensamblador o código máquina de una arquitectura concreta; en nuestro caso, como se detalla en
la sección \ref{sec:c-destino}, el lenguaje destino será \textbf{C}.

Conceptualmente, esta fase debe resolver cuestiones como:

\begin{itemize}
    \item qué instrucciones (o construcciones del lenguaje destino) representan cada operación;
    \item dónde almacenar los valores (registros, memoria, variables temporales);
    \item qué registros utilizar (\textbf{asignación de registros}), cuando el destino es
    código máquina;
    \item cómo representar las estructuras de control;
    \item cómo organizar la memoria necesaria en tiempo de ejecución.
\end{itemize}

\[
\text{IR} \rightarrow \text{selección de instrucciones} \rightarrow \text{asignación de
registros} \rightarrow \text{código destino}
\]

\subsection{Front-end y back-end}
\label{sec:frontend-backend}

Las fases del compilador pueden agruparse conceptualmente en dos grandes bloques.

El \textbf{front-end} depende principalmente del \textbf{lenguaje fuente}. Sus responsabilidades
incluyen:

\begin{itemize}
    \item análisis léxico (sección \ref{sec:lexico});
    \item análisis sintáctico (sección \ref{sec:sintactico});
    \item análisis semántico (sección \ref{sec:semantico});
    \item generación de la representación intermedia (sección \ref{sec:codigo-intermedio}).
\end{itemize}

El \textbf{back-end} recibe la representación intermedia y se ocupa de transformarla en código
del lenguaje destino, incluyendo las optimizaciones y las decisiones relacionadas con la
generación de código.

\[
\text{Fuente} \rightarrow \boxed{\text{Front-end}} \rightarrow \text{IR} \rightarrow
\boxed{\text{Back-end}} \rightarrow \text{Destino}
\]

Esta separación tiene una ventaja práctica enorme: permite que \textbf{distintos lenguajes
fuente compartan un mismo back-end},

\[
\begin{array}{ccc}
\text{Lenguaje A} & \rightarrow & \text{Front-end A} \\
\text{Lenguaje B} & \rightarrow & \text{Front-end B}
\end{array}
\rightarrow
\text{IR}
\rightarrow
\text{Back-end}
\]

o que un mismo front-end se combine con \textbf{distintos back-ends}:

\[
\text{Fuente}
\rightarrow
\text{Front-end}
\rightarrow
\text{IR}
\rightarrow
\begin{cases}
\text{Back-end x86} \\
\text{Back-end ARM}
\end{cases}
\]

En nuestro proyecto de la materia, Flex y Bison nos ayudarán a construir precisamente el
\textbf{front-end} (análisis léxico y sintáctico), mientras que nuestro back-end —más simple de
lo habitual gracias a que el destino es C— se encargará de traducir la IR a código C válido.

\subsection{Fases y pasadas}
\label{sec:fases-pasadas}

Es importante no confundir dos nociones distintas:

\begin{itemize}
    \item Las \textbf{fases} representan las distintas \textbf{tareas lógicas} que debe
    realizar el compilador (léxico, sintáctico, semántico, generación de IR, optimización,
    generación de código).
    \item Una \textbf{pasada} es una recorrida completa sobre el programa (leyendo un archivo de
    entrada y escribiendo un archivo de salida) durante la cual se ejecutan una o más fases.
\end{itemize}

Una implementación concreta puede, por ejemplo, agrupar el análisis léxico, sintáctico,
semántico y la generación de IR en una \textbf{primera pasada}:

\[
\text{Fuente} \rightarrow \text{léxico} \rightarrow \text{sintáctico} \rightarrow
\text{semántico} \rightarrow \text{IR}
\]

realizar la optimización como una \textbf{segunda pasada} (opcional):

\[
\text{IR} \rightarrow \text{IR optimizada}
\]

y finalmente generar el código destino en una \textbf{tercera pasada}:

\[
\text{IR optimizada} \rightarrow \text{código destino}
\]

En resumen: las \textbf{fases} describen \textit{qué} tareas conceptuales realiza el
compilador; las \textbf{pasadas} describen \textit{cómo} se organiza su implementación. En
nuestro compilador, es habitual que Flex y Bison colaboren dentro de una misma pasada: Bison
(el analizador sintáctico) va pidiendo tokens a Flex (el analizador léxico) \textbf{a demanda},
uno a la vez, en lugar de que Flex primero tokenice todo el archivo y luego Bison lo procese.

\subsection{Herramientas para la construcción de compiladores}
\label{sec:herramientas}

Las distintas fases de un compilador pueden implementarse manualmente o mediante herramientas
especializadas que generan automáticamente componentes a partir de especificaciones de alto
nivel. Entre las más comunes:

\begin{enumerate}
    \item \textbf{Generadores de analizadores sintácticos (parsers)}, que producen
    automáticamente analizadores sintácticos a partir de una \textbf{descripción gramatical} de
    un lenguaje de programación. \textit{Ejemplo: Bison.}
    \item \textbf{Generadores de escáneres (analizadores léxicos)}, que producen analizadores
    léxicos a partir de una descripción de los tokens de un lenguaje mediante
    \textbf{expresiones regulares}. \textit{Ejemplo: Flex.}
    \item \textbf{Motores de traducción orientados a la sintaxis}, que producen rutinas para
    recorrer el árbol sintáctico y generar código intermedio (esto es, en esencia, lo que
    hacemos al escribir acciones semánticas en las reglas de Bison).
    \item \textbf{Generadores de generadores de código}, que producen un generador de código a
    partir de reglas para traducir cada operación del lenguaje intermedio al lenguaje/arquitectura
    destino.
    \item \textbf{Motores de análisis de flujo de datos}, utilizados para obtener información
    sobre cómo se propagan los valores por el programa; son clave para muchas optimizaciones.
    \item \textbf{Kits de construcción de compiladores}, que proporcionan un conjunto integrado
    de componentes para implementar varias de estas tareas a la vez.
\end{enumerate}

Estas herramientas \textbf{no constituyen fases adicionales}: son mecanismos que facilitan la
implementación de las fases que ya existen conceptualmente. Flex y Bison son, respectivamente,
un generador de escáneres y un generador de parsers: no añaden ninguna fase nueva al esquema de
la figura de la sección \ref{sec:proceso-completo}, sólo \textbf{automatizan} la construcción de
las dos primeras.

\subsection{Proceso completo}
\label{sec:proceso-completo}

El proceso completo puede resumirse de la siguiente manera:

\[
\boxed{
\text{Fuente}
\rightarrow
\text{Análisis léxico}
\rightarrow
\text{Tokens}
}
\]

\[
\boxed{
\text{Tokens}
\rightarrow
\text{Análisis sintáctico}
\rightarrow
\text{AST}
}
\]

\[
\boxed{
\text{AST}
\rightarrow
\text{Análisis semántico}
\rightarrow
\text{AST validado}
}
\]

\[
\boxed{
\text{AST validado}
\rightarrow
\text{Código intermedio}
\rightarrow
\text{Optimización}
}
\]

\[
\boxed{
\text{IR optimizada}
\rightarrow
\text{Generación de código}
\rightarrow
\text{Código destino}
}
\]

En cada una de estas fases intervienen estructuras auxiliares transversales, como la
\textbf{tabla de símbolos} (sección \ref{sec:tabla-simbolos}) y los mecanismos de \textbf{manejo
de errores} (sección \ref{sec:errores}).

\section{Ejemplo completo: recorrido de una instrucción a través de todas las fases}
\label{sec:ejemplo-completo}

Para consolidar la relación entre todas las fases, es útil seguir un segundo ejemplo —clásico
del Libro del Dragón— algo más rico que \texttt{int x = 10 + 20;}, porque permite ver en acción
la \textbf{coerción de tipos} y una \textbf{optimización real}. Consideremos la instrucción de
asignación:

\begin{lstlisting}
posicion = inicial + velocidad * 60
\end{lstlisting}

donde \texttt{posicion}, \texttt{inicial} y \texttt{velocidad} son variables de tipo
\texttt{float}.

\paragraph{1. Análisis léxico.} Los caracteres se agrupan en lexemas y se traducen a tokens:

\[
\langle \mathbf{id}, 1\rangle\ \langle = \rangle\ \langle \mathbf{id}, 2\rangle\ \langle +
\rangle\ \langle \mathbf{id}, 3\rangle\ \langle * \rangle\ \langle 60 \rangle
\]

donde la entrada $1$ de la tabla de símbolos corresponde a \texttt{posicion}, la $2$ a
\texttt{inicial} y la $3$ a \texttt{velocidad}. El lexema \texttt{60} da lugar, en rigor, a un
token con su propia entrada en la tabla de símbolos para representar el entero.

\paragraph{2. Análisis sintáctico.} A partir de esos tokens se construye un árbol sintáctico que
refleja la precedencia usual de los operadores (la multiplicación se agrupa antes que la suma):

\begin{lstlisting}
=
|-- <id,1>            (posicion)
`-- +
    |-- <id,2>        (inicial)
    `-- *
        |-- <id,3>    (velocidad)
        `-- 60
\end{lstlisting}

El nodo raíz \texttt{=} indica que el resultado de la suma se almacena en la ubicación de
\texttt{posicion}; el nodo \texttt{+} indica que primero debe sumarse \texttt{inicial} al
resultado de la multiplicación; el nodo \texttt{*} indica que \texttt{velocidad} se multiplica
por \texttt{60} \textbf{antes} de sumarse, en concordancia con las reglas usuales de precedencia
aritmética.

\paragraph{3. Análisis semántico.} Supongamos que \texttt{posicion}, \texttt{inicial} y
\texttt{velocidad} están declaradas como \texttt{float}, mientras que el lexema \texttt{60} es,
por sí solo, un entero. El comprobador de tipos detecta que el operador \texttt{*} se aplica a
un \texttt{float} (\texttt{velocidad}) y a un \texttt{int} (\texttt{60}). Como el lenguaje
permite la \textbf{coerción} de \texttt{int} a \texttt{float} en operaciones aritméticas
mixtas, el analizador semántico inserta de manera explícita un operador de conversión
\texttt{inttofloat} sobre el \texttt{60}. El árbol resultante incorpora entonces ese nodo
adicional antes de pasar a la siguiente fase.

\paragraph{4. Generación de código intermedio.} Usando código de tres direcciones, con una
instrucción por cada operador del árbol (incluyendo la conversión de tipos), se obtiene:

\begin{lstlisting}
t1 = inttofloat(60)
t2 = id3 * t1
t3 = id2 + t2
id1 = t3
\end{lstlisting}

Aquí \texttt{id1}, \texttt{id2} e \texttt{id3} son las referencias a las entradas de la tabla de
símbolos de \texttt{posicion}, \texttt{inicial} y \texttt{velocidad}, respectivamente, y
\texttt{t1}, \texttt{t2}, \texttt{t3} son nombres temporales generados por el compilador.
Obsérvese que cada instrucción de tres direcciones tiene, como máximo, un operador del lado
derecho, lo cual fuerza a explicitar el orden en el que se realizan las operaciones (la
multiplicación queda calculada antes que la suma, tal como exigía el árbol sintáctico).

\paragraph{5. Optimización de código.} El optimizador puede advertir que la conversión de
\texttt{60} a \texttt{float} puede resolverse en tiempo de compilación, reemplazando el entero
\texttt{60} por la constante \texttt{60.0}; también puede notar que \texttt{t3} se usa una
única vez, así que su valor puede propagarse directamente a \texttt{id1}. El resultado es una
secuencia más corta:

\begin{lstlisting}
t1 = id3 * 60.0
id1 = id2 + t1
\end{lstlisting}

\paragraph{6. Generación de código destino.} Si el destino fuera código máquina, y se usaran los
registros \texttt{R1} y \texttt{R2}, podría generarse:

\begin{lstlisting}
LDF  R2, id3
MULF R2, R2, #60.0
LDF  R1, id2
ADDF R1, R1, R2
STF  id1, R1
\end{lstlisting}

En nuestro proyecto, en cambio, el back-end no genera ensamblador ni código máquina: traduce la
IR optimizada a \textbf{código C}, delegando la selección de instrucciones y la asignación de
registros al compilador de C. Para la secuencia anterior, el back-end podría emitir algo
equivalente a:

\begin{lstlisting}
float t1;
t1 = velocidad * 60.0;
posicion = inicial + t1;
\end{lstlisting}

Este recorrido completo —léxico, sintáctico, semántico (con coerción), IR, optimización y
generación de código— es el mismo esquema que aplicaremos, en miniatura, a
\texttt{int x = 10 + 20;} durante la práctica con Flex y Bison.

\section{Uso de C como lenguaje destino}
\label{sec:c-destino}

En nuestro compilador se utilizará \textbf{C} como lenguaje destino. Esto significa que el
compilador \textbf{no generará} directamente código máquina ni ensamblador: en su lugar,
generará un programa escrito en C que representa las operaciones del lenguaje fuente. El
proceso completo será:

\[
\begin{array}{c}
\text{Lenguaje fuente}
\rightarrow
\text{Análisis léxico}
\rightarrow
\text{Análisis sintáctico}
\rightarrow
\text{Análisis semántico}
\\[6pt]
\downarrow \\[6pt]
\text{IR}
\rightarrow
\text{Código C}
\rightarrow
\text{Compilador de C}
\rightarrow
\text{Código máquina}
\end{array}
\]

El objetivo de nuestro back-end será, por lo tanto, traducir la representación intermedia a
código C válido. Por ejemplo, si la IR contiene:

\begin{lstlisting}
t1 = 10 + 20
x = t1
\end{lstlisting}

el back-end podría generar:

\begin{lstlisting}
int x;
x = 10 + 20;
\end{lstlisting}

y será el compilador de C ya existente quien se encargue de las etapas necesarias para producir
código específico de la plataforma final.

Esta decisión simplifica considerablemente la construcción del back-end: los compiladores de C
disponibles ofrecen optimizaciones maduras y soporte para múltiples arquitecturas, por lo que
nuestro proyecto no necesita reimplementar instrucciones, registros ni convenciones específicas
de cada arquitectura (lo que sí sería necesario si generáramos código máquina directamente).
La utilización de C como lenguaje destino permite concentrar el esfuerzo del proyecto en el
\textbf{diseño y procesamiento de nuestro propio lenguaje} —justamente, el front-end que
construiremos con Flex y Bison— delegando en el compilador de C la generación final de código
para la plataforma correspondiente.

\section{De la teoría del Capítulo 1 a la práctica: bases para Flex y Bison}
\label{sec:flex-bison}

Esta sección conecta explícitamente los conceptos ya vistos con las herramientas que
comenzaremos a usar mañana.

\subsection{Expresiones regulares como patrones léxicos: la base de Flex}

En la sección \ref{sec:lexemas-tokens-patrones} definimos un \textbf{patrón} como la regla que
describe qué cadenas de caracteres pueden formar los lexemas de un token dado. El modelo formal
estándar para expresar esas reglas de manera precisa y compacta son las \textbf{expresiones
regulares}, apoyadas en el modelo de las \textbf{máquinas de estados finitos} (autómatas
finitos). Estos son, junto con las gramáticas libres de contexto, de los modelos matemáticos
más básicos usados en el diseño de compiladores.

\textbf{Flex} es un \textbf{generador de analizadores léxicos}: recibe un archivo de
especificación (típicamente \texttt{.l} o \texttt{.flex}) formado por una lista de pares
\textit{(expresión regular, acción en C)}, y genera automáticamente el código C de una función
(\texttt{yylex()}) que:

\begin{enumerate}
    \item lee el flujo de caracteres del programa fuente;
    \item en cada paso, encuentra el \textbf{lexema más largo} que coincide con alguno de los
    patrones especificados (la llamada regla de la \textit{coincidencia más larga});
    \item ejecuta la acción en C asociada, que típicamente construye y devuelve el
    \textbf{token} correspondiente —exactamente el par $\langle \text{nombre-token},
    \text{valor-atributo} \rangle$ descrito en la sección \ref{sec:lexico}—;
    \item ignora (o trata especialmente) los espacios en blanco, tal como hace cualquier
    analizador léxico.
\end{enumerate}

En otras palabras: \textbf{todo lo que hicimos ``a mano'' en la sección \ref{sec:lexico}
(agrupar caracteres en lexemas, asignarles una categoría) es exactamente lo que Flex
automatiza}, a partir de una tabla de patrones como la de la tabla \ref{tab:tokens}.

\subsection{Gramáticas libres de contexto y árboles sintácticos: la base de Bison}

En la sección \ref{sec:especificacion-sintactica} vimos que la sintaxis de un lenguaje se
especifica mediante una \textbf{gramática libre de contexto}: un conjunto de producciones que
indican cómo se combinan los tokens (símbolos terminales) y las categorías sintácticas
(símbolos no terminales) para formar construcciones válidas.

\textbf{Bison} es un \textbf{generador de analizadores sintácticos}: recibe un archivo de
especificación (típicamente \texttt{.y} o \texttt{.yacc}) formado por una gramática —muy similar
en notación a la de la sección \ref{sec:especificacion-sintactica}— junto con \textbf{acciones
semánticas} en C asociadas a cada producción, y genera automáticamente el código C de una
función (\texttt{yyparse()}) capaz de:

\begin{enumerate}
    \item pedir tokens, uno a la vez, al analizador léxico generado por Flex (\texttt{yylex()});
    \item determinar si la secuencia de tokens recibida es derivable desde el símbolo inicial de
    la gramática, usando un algoritmo de análisis ascendente de tipo \textbf{LALR};
    \item ejecutar, en cada reducción de una producción, la acción semántica en C asociada —lo
    que en la sección \ref{sec:ast} llamamos \textbf{traducción dirigida por la sintaxis}—,
    típicamente para construir un nodo del \textbf{árbol de sintaxis abstracta};
    \item reportar un error sintáctico (invocando \texttt{yyerror}) cuando ninguna producción de
    la gramática permite continuar el análisis.
\end{enumerate}

Es decir: \textbf{todo lo que hicimos ``a mano'' en las secciones \ref{sec:sintactico} y
\ref{sec:ast} (construir el árbol sintáctico y luego el AST a partir de la gramática de la
sección \ref{sec:especificacion-sintactica}) es exactamente lo que Bison automatiza}.

\subsection{El flujo Flex + Bison y su correspondencia con las fases del compilador}

En una implementación típica con estas herramientas, Flex y Bison \textbf{no se ejecutan como
dos pasadas separadas}: Bison controla el proceso y pide tokens a Flex \textbf{bajo demanda},
uno por vez, tal como se explicó en la sección \ref{sec:fases-pasadas}. La correspondencia con
las fases del capítulo es directa:

\begin{table}[H]
\centering
\begin{tabular}{lll}
\toprule
\textbf{Fase del compilador} & \textbf{Herramienta} & \textbf{Entrada / salida} \\
\midrule
Análisis léxico & Flex (\texttt{yylex}) & caracteres $\to$ tokens \\
Análisis sintáctico & Bison (\texttt{yyparse}) & tokens $\to$ árbol / AST \\
Traducción dirigida por la sintaxis & acciones en Bison & AST anotado \\
Análisis semántico & código C del programador & AST validado \\
\bottomrule
\end{tabular}
\caption{Correspondencia entre las fases del Capítulo 1 y el par Flex/Bison}
\label{tab:flex-bison-fases}
\end{table}

En síntesis: \textbf{Flex y Bison automatizan el front-end} del compilador (secciones
\ref{sec:lexico} a \ref{sec:codigo-intermedio}), pero \textbf{no reemplazan} el diseño
conceptual que vimos en este capítulo: antes de escribir una sola línea de \texttt{.l} o
\texttt{.y}, necesitamos tener claros los \textbf{patrones léxicos} (expresiones regulares) y la
\textbf{gramática} de nuestro lenguaje, exactamente como se plantearon en las tablas y
producciones de las secciones \ref{sec:lexemas-tokens-patrones} y
\ref{sec:especificacion-sintactica}. Ese es, precisamente, el contenido que repasaremos en la
clase de mañana antes de comenzar con la implementación.

\section{Panorama complementario}

\subsection{La evolución de los lenguajes de programación}

Las primeras computadoras (década de 1940) se programaban directamente en \textbf{lenguaje
máquina}. El desarrollo de los \textbf{lenguajes ensambladores}, a inicios de 1950, introdujo
mnemónicos y, más adelante, macroinstrucciones. En la segunda mitad de esa década surgieron los
primeros \textbf{lenguajes de alto nivel} (Fortran, Cobol, Lisp), pensados para que la
programación fuera más natural en distintos dominios.

Una clasificación habitual es por \textbf{generación}: 1\textsuperscript{a} generación
(lenguaje máquina), 2\textsuperscript{a} (ensambladores), 3\textsuperscript{a} (lenguajes de
alto nivel de propósito general, como C, C++, Java), 4\textsuperscript{a} (lenguajes de
aplicación específica, como SQL) y 5\textsuperscript{a} (lenguajes basados en lógica y
restricciones, como Prolog). Otra clasificación distingue lenguajes \textbf{imperativos}
(especifican \textit{cómo} realizar un cálculo, con una noción explícita de estado, como C) de
los \textbf{declarativos} (especifican \textit{qué} calcular, como los lenguajes funcionales o
de lógica).

Cada avance en los lenguajes de programación impuso nuevas exigencias sobre quienes escriben
compiladores, que debieron idear algoritmos y representaciones para dar soporte a las nuevas
características, aprovechando además la evolución de la arquitectura de las máquinas.

\subsection{La ciencia de construir un compilador}

El diseño de compiladores combina, de forma ejemplar, \textbf{teoría matemática} y
\textbf{validación empírica}: se formula una abstracción matemática que capture las
características clave de un problema (por ejemplo, autómatas finitos y expresiones regulares
para el léxico, gramáticas libres de contexto para la sintaxis, árboles para representar
programas) y se resuelve con técnicas matemáticas, refinando la solución con experimentación.

Toda optimización de código debe cumplir ciertos objetivos de diseño, en orden de importancia:

\begin{enumerate}
    \item Debe ser \textbf{correcta}: preservar el significado del programa compilado. Este es,
    con diferencia, el objetivo más importante: un compilador que genere código incorrecto no
    sirve, sin importar qué tan rápido sea el código resultante.
    \item Debe \textbf{mejorar el rendimiento} de una amplia variedad de programas de entrada.
    \item El \textbf{tiempo de compilación} debe mantenerse razonable, para no entorpecer el
    ciclo de desarrollo y depuración.
    \item El \textbf{esfuerzo de ingeniería} requerido debe ser manejable.
\end{enumerate}

Vale la pena notar que muchos de los problemas centrales de la optimización de compiladores son
\textbf{indecidibles} en general; por eso buena parte del trabajo consiste en formular
heurísticas razonables, más que en buscar soluciones perfectas.

\subsection{Aplicaciones de la tecnología de compiladores}

Las técnicas desarrolladas para construir compiladores se aplican mucho más allá de la
compilación de lenguajes de programación de propósito general. Algunos ejemplos mencionados en
el capítulo:

\begin{itemize}
    \item \textbf{Traducción binaria}: convertir código máquina de una arquitectura a otra.
    \item \textbf{Síntesis de hardware}: traducir descripciones en lenguajes como VHDL o Verilog
    a circuitos concretos.
    \item \textbf{Intérpretes de consultas de bases de datos}: compilar o interpretar lenguajes
    como SQL.
    \item \textbf{Simulación compilada}: compilar un diseño para simularlo mucho más rápido que
    interpretándolo.
    \item \textbf{Herramientas de productividad de software}: comprobación de tipos, análisis de
    flujo de datos para detectar errores (referencias nulas, desbordamientos de buffer) y
    herramientas de administración de memoria (por ejemplo, recolección de basura), todas ellas
    apoyadas en técnicas de análisis de programas nacidas en el contexto de la optimización de
    compiladores.
\end{itemize}

\section{Conclusión}

A partir del modelo presentado en este capítulo, el desarrollo de nuestro compilador puede
organizarse como una secuencia de etapas claramente diferenciadas.

En primer lugar, debemos definir la \textbf{especificación léxica} del lenguaje: sus tokens,
lexemas y patrones (expresiones regulares), que es exactamente lo que especificaremos en el
archivo de Flex. Luego debemos establecer la \textbf{gramática} que determina cómo se combinan
esos tokens, que es lo que especificaremos en el archivo de Bison, junto con las acciones
semánticas que construyen el AST. Sobre esa estructura incorporaremos las \textbf{reglas
semánticas} y la \textbf{tabla de símbolos} necesarias para validar los programas. Una vez
definido el front-end, estableceremos una \textbf{representación intermedia} que exprese las
operaciones del lenguaje de manera independiente del código C generado, y finalmente
implementaremos el \textbf{back-end} encargado de traducir esa representación a C.

De esta manera, el desarrollo seguirá la organización clásica de un compilador:

\[
\boxed{
\text{Lenguaje}
\rightarrow
\text{Léxico}
\rightarrow
\text{Sintaxis}
\rightarrow
\text{Semántica}
\rightarrow
\text{IR}
\rightarrow
\text{C}
}
\]

y, en la práctica concreta de la cátedra, las dos primeras flechas de ese esquema
—\textbf{léxico} y \textbf{sintaxis}— son precisamente el territorio de \textbf{Flex} y
\textbf{Bison}, las herramientas que comenzaremos a estudiar en la próxima clase.

\end{document}
